\documentclass{situationist-preprint}

\title{The Irrationality of Procrastination}
\msid{ACT-13-a}
\msfamily{ACT-13}
\msdiscipline{Action, Agency and Free Will}
\mstype{Principal manuscript}
\mskeywords{procrastination; akrasia; temporal discounting; hyperbolic discounting; preference reversal; self-control}

\begin{document}
\maketitle

\begin{abstract}
Procrastination is common, and those who engage in it usually regret it. Yet it is not obvious that it is irrational: delay is often reasonable, and a person who postpones an unpleasant task may simply weigh the present more than the future, which some theories of rational choice permit. We argue that procrastination is irrational, because it is a form of akrasia involving temporal discounting. The procrastinator delays an intended action while judging, at the time, that delaying is worse than acting now; she therefore acts against her better judgment. What explains the gap between her judgment and her action is that the motivational pull of an option, unlike its assessed value, is discounted steeply with temporal distance, so that the present costs of the task loom larger than its future benefits at every moment at which she could begin. We distinguish procrastination from reasonable delay, show that the discounting involved produces preference reversals that the agent herself foresees and rejects, and answer objections from the view that procrastination involves no single irrational choice.
\end{abstract}

\section{The problem}

The economist who postpones a trivial errand for eight months, and the student who begins an essay the night before it is due, are familiar figures \citep{Akerlof1991}. In the psychological literature, procrastination is defined as voluntarily delaying an intended course of action despite expecting to be worse off for the delay \citep{Steel2007}. On this definition, the question whether procrastination is irrational seems to answer itself. But two lines of thought resist the answer. On the first, delay may express a rational time preference: if an agent may discount future welfare, as many models of choice allow, then preferring a smaller present comfort to a larger future benefit is a permissible preference, not a failure of reason. On the second, procrastination does not consist in any single choice that is worse than its alternative: each day's delay makes no difference, and only the series is bad \citep{Andreou2007}. Economists have modelled procrastination as the product of present bias, a discount function that weights the present disproportionately relative to all later times \citep{ODonoghueRabin1999}; and the essays collected by \citet{AndreouWhite2010} examine its relation to weakness of will. We argue, consistently with ACT-04-a, that procrastination is a form of akrasia, and that temporal discounting explains why it takes the form it does.

\section{Framework}

\begin{definition}[Procrastination]\label{d:proc}
An agent \emph{procrastinates} with respect to a task if she intends to perform it, and at a time at which she can begin it she voluntarily delays it, while judging that beginning now would be better, all things considered, than delaying.
\end{definition}

\begin{definition}[Discounting]\label{d:disc}
The \emph{motivational discount} of an option is the reduction in its motivational strength, in the sense of ACT-04-a, as its costs or benefits recede in time. Discounting is \emph{hyperbolic} if the discount rate is steep for short delays and shallow for long ones, so that the relative pull of two options can reverse as they approach.
\end{definition}

\begin{principle}[Value is not discounted by proximity]\label{p:value}
An agent's all-things-considered judgment of the value of an option does not vary with its temporal proximity, except where proximity changes the facts that determine the value.
\end{principle}

Definition~\ref{d:proc} follows the psychological definition, with the agent's expectation of being worse off made explicit as a better judgment held at the time. It excludes reasonable delay: the agent who postpones a task because she judges that she will do it better later, or because something more important has arisen, does not procrastinate. Principle~\ref{p:value} is the observation on which ACT-04-a relies: the agent who judges in the afternoon that she should start the essay today judges the same in the evening, when what has changed is only that the costs of starting are now immediate. A judgment of value that tracked proximity as such would revise itself merely because time had passed, which is not a response to any reason. The empirical evidence for hyperbolic discounting of motivation is extensive \citep{Ainslie1992}; its consequence for dynamic choice, that the agent's later choices depart from her earlier plans, was identified by \citet{Strotz1955}.

\section{The argument}

\begin{lemma}\label{l:akr}
Every case of procrastination is a case of strict akrasia.
\end{lemma}

\begin{proof}
By Definition~\ref{d:proc}, the procrastinator voluntarily delays while judging that beginning now would be better, all things considered. Delaying is something she does intentionally, for a reason she recognizes, the avoidance of present unpleasantness. Her action is free provided she could have begun, which is the case if her mechanism is responsive to reasons, as it is in the normal procrastinator, who begins when the deadline makes the costs of delay vivid. She therefore intentionally and freely performs an option other than the one her better judgment favors, while holding that judgment. This is strict akrasia as ACT-04-a defines it.
\end{proof}

\begin{lemma}\label{l:disc}
In procrastination, the strength of the motive to delay exceeds that of the motive supplied by the better judgment because of temporal discounting.
\end{lemma}

\begin{proof}
The task's costs (effort, boredom, anxiety) fall at the time of beginning; its benefits (completion, avoiding the costs of lateness) fall later. At a time well before the possible start, both are distant, and by hyperbolic discounting (Definition~\ref{d:disc}) their motivational weights stand in roughly the proportion of their values, so the agent both judges and wants to begin on time; this is why she forms the intention. As the start approaches, the costs become immediate and their motivational weight rises steeply, while the benefits remain distant. The relative pull reverses. By Principle~\ref{p:value}, her judgment does not reverse with it. The gap between judgment and motivation that Lemma~\ref{l:akr} identifies is thus produced by discounting.
\end{proof}

\begin{theorem}\label{t:main}
Procrastination is irrational: it is a form of akrasia involving temporal discounting.
\end{theorem}

\begin{proof}
By Lemma~\ref{l:akr}, procrastination is akrasia, and by Lemma~\ref{l:disc} the akrasia is produced by temporal discounting. ACT-04-a concludes that akrasia is a failure of rationality: the akratic agent fails to be moved in accordance with a judgment she holds, by a desire she herself ranks lower. So procrastination is irrational. The discounting adds a second, diachronic, defect: the agent's motivations at successive times rank the same options in opposite orders, which she foresees when she forms her intention and rejects when she acts against it.
\end{proof}

\section{Objections and replies}

\begin{objection}[Rational time preference]
Pure time preference, a preference for goods sooner rather than later, is not irrational. The procrastinator merely gives the present the weight she prefers.
\end{objection}

\begin{reply}
Whether a constant rate of pure time preference is rational is a question we need not settle. The procrastinator's discounting is not of that kind: it is hyperbolic, producing reversals that her own judgment, formed at an earlier time and retained at the time of action, rejects. If she endorsed the steep discount, she would judge delay better and would not procrastinate in the sense of Definition~\ref{d:proc}; she would delay reasonably by her own lights. The typical procrastinator does not endorse it, as her regret and her use of deadlines and commitment devices show \citep{Ainslie1992}. The companion manuscript ACT-13-b examines deadlines as instruments of self-control.
\end{reply}

\begin{objection}[No single choice]
Each day's delay is harmless, since the task can be done tomorrow as well as today. Only the series of delays is bad. The procrastinator therefore makes no choice against her better judgment on any occasion \citep{Andreou2007}.
\end{objection}

\begin{reply}
Andreou's cases show that procrastination can be a failure to settle when to act, in which no single delay is worse than its alternative. But in such cases the agent's better judgment is about policy: she judges that she should begin now rather than continue a pattern of delay she knows to be bad, since nothing distinguishes today from tomorrow except that today is the day she is deciding. Acting against that judgment is akratic. Where she judges, at each moment, that delaying one more day is genuinely as good as starting, she is not procrastinating by Definition~\ref{d:proc}, though she may be making an error about vague thresholds of the kind Andreou describes.
\end{reply}

\begin{objection}[Revised judgment]
Many procrastinators do not act against their judgment; they talk themselves into thinking that the task can wait. This is unreasonable revision of a resolution, not akrasia.
\end{objection}

\begin{reply}
We agree that such cases occur. If the revision is genuine, the agent is not procrastinating under Definition~\ref{d:proc} at the moment of delay, though she did so earlier if she delayed while still judging that she should not. Self-deceived revision of this kind is itself a product of the discounting that Lemma~\ref{l:disc} describes, and it is also irrational, by Principle~\ref{p:value}, since the revision responds to proximity rather than to reasons.
\end{reply}

\section{Conclusion}

The procrastinator delays what she judges, at the time, she should do now. She is moved by the immediacy of the task's costs, whose motivational weight rises steeply as they approach, while her judgment of value stays where it was. Procrastination is therefore a form of akrasia involving temporal discounting, and it is irrational for the reason that akrasia is.

\begin{thebibliography}{99}\small
\bibitem[Ainslie(1992)]{Ainslie1992} Ainslie, G. (1992). \emph{Picoeconomics: The Strategic Interaction of Successive Motivational States within the Person}. Cambridge University Press.
\bibitem[Akerlof(1991)]{Akerlof1991} Akerlof, G.~A. (1991). Procrastination and obedience. \emph{American Economic Review}, 81(2), 1--19.
\bibitem[Andreou(2007)]{Andreou2007} Andreou, C. (2007). Understanding procrastination. \emph{Journal for the Theory of Social Behaviour}, 37(2), 183--193.
\bibitem[O'Donoghue and Rabin(1999)]{ODonoghueRabin1999} O'Donoghue, T. and Rabin, M. (1999). Doing it now or later. \emph{American Economic Review}, 89(1), 103--124.
\bibitem[Steel(2007)]{Steel2007} Steel, P. (2007). The nature of procrastination: a meta-analytic and theoretical review of quintessential self-regulatory failure. \emph{Psychological Bulletin}, 133(1), 65--94.
\bibitem[Strotz(1955)]{Strotz1955} Strotz, R.~H. (1955). Myopia and inconsistency in dynamic utility maximization. \emph{The Review of Economic Studies}, 23(3), 165--180.
\bibitem[Andreou and White(2010)]{AndreouWhite2010} Andreou, C. and White, M.~D. (eds.) (2010). \emph{The Thief of Time: Philosophical Essays on Procrastination}. Oxford University Press.
\end{thebibliography}

\end{document}
